\subsubsection{Earlier Sequential Application Control}
\label{sec:joint}
\paragraph{Publication and answer times.}
All 12 application episodes complete. Current answers pass their supported
checks and retrieval locates the new source in v2 in every episode. Publication and retrieval therefore work in these episodes. We still need
the answer checks to determine whether the new outputs are valid.

\begin{table}[t]
\caption{Median times from update arrival across three matched episode blocks
with different new documents. $T_c$ is current completion, $T_p$ is publication and $T_f$ is new-answer completion. These are descriptive medians, not confidence bounds.}
\label{tab:joint}
\centering
\begin{tabular}{lrrr}
\toprule
Policy & $T_c$ & $T_p$ & $T_f$\\
\midrule
Standard concurrent & 76.2 & 14.7 & 189.8 \\
Standard deferred & 70.7 & 76.0 & 189.0 \\
Full prototype & 78.5 & 10.3 & 198.9 \\
Skip4 prototype & 76.6 & 9.9 & 189.6 \\
\bottomrule

\end{tabular}
\end{table}
The full prototype publishes earlier than standard concurrent (10.3 versus
14.7\,s median), but its current and fresh answers finish later (78.5 versus
76.2\,s for current answers and 198.9 versus 189.8\,s for new answers). Skip4 publishes at 9.9\,s, with fresh completion
189.6\,s, essentially the same scale as standard. It is already ineligible under
the separate-document paired gate. Deferral protects current service (70.7\,s)
but delays publication to 76.0\,s. None satisfies the frozen automatic joint
benefit criterion.

Figure~\ref{fig:research-episode} shows the timing of one measured episode. In every
concurrent episode, publication precedes current-answer completion. Under the
frozen single-slot trace, Equation~\ref{eq:critical-path} is then governed by
current inference and the next request rather than publication. The measured result is that publication is already complete while
current inference remains active, so accelerating publication does not advance
the next request in these episodes. It does not
prove that publication speed is irrelevant for earlier-arriving independent
requests, multiple slots or another workload.
\begin{figure*}[t]
\centering
\begin{tikzpicture}[font=\scriptsize,>=Latex]
\fill[answerblue] (0,.72) rectangle ++(.25,.16);\node[anchor=west] at (.33,.80) {Current inference};
\fill[updatepurple] (3.1,.72) rectangle ++(.25,.16);\node[anchor=west] at (3.43,.80) {Incremental publication};
\fill[freshorange] (7.0,.72) rectangle ++(.25,.16);\node[anchor=west] at (7.33,.80) {New-version inference};
\node[anchor=east,font=\scriptsize] at (-.18,-0.120) {Standard concurrent};
\fill[answerblue] (0,0.000) rectangle (4.1315,0.200);
\fill[updatepurple] (0.0000,-0.250) rectangle (0.7919,-0.050);
\draw[black!45,densely dotted] (4.1315,0.100) -- (4.1402,0.100);
\fill[freshorange] (4.1402,0.000) rectangle (10.2518,0.200);
\node[anchor=west,font=\scriptsize] at (10.3518,0.100) {189.8s};
\draw[updatepurple] (0.7919,-0.300) -- (0.7919,-0.460);
\node[anchor=north,font=\scriptsize,text=updatepurple] at (0.7919,-0.460) {14.7s};
\node[anchor=east,font=\scriptsize] at (-.18,-1.240) {Standard deferred};
\fill[answerblue] (0,-1.120) rectangle (3.8336,-0.920);
\fill[updatepurple] (3.8336,-1.370) rectangle (4.1127,-1.170);
\draw[black!45,densely dotted] (4.1127,-1.020) -- (4.1182,-1.020);
\fill[freshorange] (4.1182,-1.120) rectangle (10.2037,-0.920);
\node[anchor=west,font=\scriptsize] at (10.3037,-1.020) {189.0s};
\draw[updatepurple] (4.1127,-1.420) -- (4.1127,-1.580);
\node[anchor=north,font=\scriptsize,text=updatepurple] at (4.1127,-1.580) {76.2s};
\node[anchor=east,font=\scriptsize] at (-.18,-2.360) {Full prototype};
\fill[answerblue] (0,-2.240) rectangle (4.2989,-2.040);
\fill[updatepurple] (0.0000,-2.490) rectangle (0.5554,-2.290);
\draw[black!45,densely dotted] (4.2989,-2.140) -- (4.3068,-2.140);
\fill[freshorange] (4.3068,-2.240) rectangle (10.7381,-2.040);
\node[anchor=west,font=\scriptsize] at (10.8381,-2.140) {198.9s};
\draw[updatepurple] (0.5554,-2.540) -- (0.5554,-2.700);
\node[anchor=north,font=\scriptsize,text=updatepurple] at (0.5554,-2.700) {10.3s};
\node[anchor=west,font=\scriptsize,text=red!70!black] at (4.4068,-2.530) {correct fact; citation-format failure};
\node[anchor=east,font=\scriptsize] at (-.18,-3.480) {Skip4 prototype};
\fill[answerblue] (0,-3.360) rectangle (4.1341,-3.160);
\fill[updatepurple] (0.0000,-3.610) rectangle (0.5432,-3.410);
\draw[black!45,densely dotted] (4.1341,-3.260) -- (4.1429,-3.260);
\fill[freshorange] (4.1429,-3.360) rectangle (10.2359,-3.160);
\node[anchor=west,font=\scriptsize] at (10.3359,-3.260) {189.6s};
\draw[updatepurple] (0.5432,-3.660) -- (0.5432,-3.820);
\node[anchor=north,font=\scriptsize,text=updatepurple] at (0.5432,-3.820) {10.1s};

\draw[->,black!60] (0,-4.28)--(11.15,-4.28);
\foreach \t in {0,50,100,150,200}{\draw[black!60] (\t*.054,-4.22)--(\t*.054,-4.34);\node[anchor=north] at (\t*.054,-4.37) {\t};}
\node at (5.4,-4.92) {Elapsed seconds since update arrival};
\end{tikzpicture}
\caption{Measured critical path, first predeclared matched episode (Normandy),
not a timeline assembled from unrelated medians. Purple ticks label publication;
right labels mark new-answer completion. Publication precedes current-answer
completion for all concurrent arms, so accelerating publication does not release
the serial inference slot earlier. Deferral instead delays publication until the
current answer, leaving fresh completion similar in this episode. The full
prototype answers ``France'' but emits an invalid citation prefix; its timing
is preserved as a failed task. Dotted gaps are post-gate retrieval/submission.}
\label{fig:research-episode}
\end{figure*}


\paragraph{Answer checks.}
The first two added documents ask about Normandy's country and Southern
California's abbreviation. Standard answers ``France'' and ``SoCal'' with valid
locators. Full prototype's first answer has the correct fact but an extra
\texttt{source\_id:} prefix in the locator, which fails the contract. The third
asks what kind of economy Victoria has. All arms generate a relevant elaboration
containing the gold ``highly diversified'', but fail exact-reference scoring.
The longer answer fails the scoring rule. Human review is needed before
interpreting it as an error in the economic description.

The original all-episode quality gate fails for standard as well as the
candidates. We retain that failed result rather than relax the rule after reading the outputs. Separately, even a more permissive interpretation of that
shared elaboration would not make full prototype dominate the measured latency
vector: it publishes faster while answering later. Timing and answer validity are separate results. A new semantic assessment could change quality labels
but would not change these measured timestamps or establish statistical benefit.

The current response uses old pinned oracle evidence, whereas the new response
uses actual retrieval. We therefore claim a versioned-update prototype with one
retrieved fresh question per episode, not a fully validated production RAG
trace. The workload demonstrates neither domain-specific household deployment
utility nor a real-time service contract. Observed response times are in minutes.
